\begin{helper}
Let $V$ be finite, let $p\in[0,1]$, and let $\nu$ be a measure on
$\mathcal P(V)\times\mathbb R^V$ with the discrete and Borel product
measurable structure. Suppose for each $F\subseteq V$,
\[
 \nu\{\eta:\eta_{\rm act}=F\}
 =\operatorname{ofReal}\bigl(p^{|F|}(1-p)^{|V|-|F|}\bigr).
\]
For any $A\subseteq S\subseteq V$,
\[
 \nu\{\eta:\eta_{\rm act}\cap S=A\}
 =\operatorname{ofReal}\bigl(p^{|A|}(1-p)^{|S|-|A|}\bigr).
\]
\end{helper}
\begin{proof}
Write $C=V\setminus S$. The activation sets with trace $A$ on $S$
are precisely $F=A\cup E$ for $E\subseteq C$; the inverse correspondence
is $E=F\setminus S$. Their activation fibres are measurable and pairwise
disjoint. Finite additivity and the assumed fibre masses therefore give
the real sum of $p^{|A\cup E|}(1-p)^{|V|-|A\cup E|}$, converted by
$\operatorname{ofReal}$. This conversion commutes with the finite sum
because $p$ and $1-p$ are nonnegative.
Since $A\cap E=\varnothing$, we have
$|A\cup E|=|A|+|E|$ and
$|V|-|A\cup E|=(|S|-|A|)+(|C|-|E|)$.
Factoring the common powers leaves
\[
 p^{|A|}(1-p)^{|S|-|A|}
 \sum_{E\subseteq C}p^{|E|}(1-p)^{|C|-|E|}.
\]
To evaluate the remaining sum, expand
$\prod_{v\in C}(p+(1-p))$ by choosing the first summand exactly at the
indices in $E$. The term for $E$ is exactly
$p^{|E|}(1-p)^{|C|-|E|}$. Thus the sum is
$(p+(1-p))^{|C|}=1$. This proves the formula. The calculation involves
only sums and products, never division, and so includes $p=0$, $p=1$,
and empty $A$, $S$, or $C$, with zeroth powers equal to one.
\end{proof}
